\chapter{General Context and Existing Study}

\section{Overview of Messaging Systems}

Instant messaging systems have evolved significantly since the early days of the internet. Modern messaging platforms can be categorized into three main architectural models:

\subsection{Centralized Messaging Systems}

Centralized systems like WhatsApp, Telegram, and Facebook Messenger operate with a single authority controlling all infrastructure. Messages flow through central servers that have technical capability to access content. While some platforms implement end-to-end encryption, the centralized architecture creates inherent trust dependencies.

\textbf{Advantages:}
\begin{itemize}
    \item Simple user experience with centralized account management
    \item Efficient message delivery and synchronization
    \item Easy implementation of features like group chats and media sharing
    \item Centralized moderation and spam prevention
\end{itemize}

\textbf{Disadvantages:}
\begin{itemize}
    \item Single point of failure for privacy and security
    \item Metadata collection (who communicates with whom, when, how often)
    \item Vulnerability to government surveillance and legal requests
    \item Risk of data breaches exposing user information
\end{itemize}

\subsection{Federated Messaging Systems}

Federated systems like Matrix and XMPP distribute control across multiple servers operated by different entities. Users on different servers can communicate through standardized protocols.

\textbf{Advantages:}
\begin{itemize}
    \item No single point of control or failure
    \item Users can choose their service provider
    \item Open protocols enable interoperability
    \item Reduced metadata exposure to any single entity
\end{itemize}

\textbf{Disadvantages:}
\begin{itemize}
    \item Complex setup and configuration for end users
    \item Metadata still visible to federated servers
    \item Inconsistent security implementations across servers
    \item Difficulty in implementing advanced features consistently
\end{itemize}

\subsection{Peer-to-Peer Messaging Systems}

Peer-to-peer systems like Briar and Tox eliminate central servers entirely, with messages transmitted directly between users or through distributed networks.

\textbf{Advantages:}
\begin{itemize}
    \item Maximum privacy with no central authority
    \item Resistance to censorship and surveillance
    \item No metadata collection by service providers
    \item True decentralization
\end{itemize}

\textbf{Disadvantages:}
\begin{itemize}
    \item Requires both parties to be online simultaneously
    \item Complex NAT traversal and connectivity issues
    \item Limited scalability for group communications
    \item Difficult user experience for non-technical users
\end{itemize}

\section{Privacy and Security Challenges}

\subsection{Metadata Exposure}

Even with end-to-end encryption, messaging systems expose significant metadata:
\begin{itemize}
    \item Communication patterns (who talks to whom)
    \item Timing information (when messages are sent)
    \item Message frequency and volume
    \item User location and IP addresses
    \item Device information and online status
\end{itemize}

This metadata can reveal sensitive information about users' relationships, habits, and activities, even when message content remains encrypted.

\subsection{Key Management Challenges}

Secure messaging requires solving complex key management problems:
\begin{itemize}
    \item \textbf{Key Generation:} Ensuring cryptographically secure random key generation
    \item \textbf{Key Storage:} Protecting private keys from unauthorized access
    \item \textbf{Key Backup:} Enabling key recovery without compromising security
    \item \textbf{Key Rotation:} Periodically updating keys to limit exposure
    \item \textbf{Key Verification:} Allowing users to verify correspondent identities
\end{itemize}

\subsection{Authentication vs. Privacy Trade-offs}

Traditional authentication systems conflict with privacy goals:
\begin{itemize}
    \item Phone number requirements link identities to real-world persons
    \item Email verification creates traceable digital footprints
    \item Username systems enable enumeration attacks
    \item Password recovery mechanisms can bypass encryption
\end{itemize}

\section{Limitations of Centralized Messaging Platforms}

\subsection{Trust Dependencies}

Centralized platforms require users to trust that:
\begin{itemize}
    \item The service provider implements encryption correctly
    \item Server operators do not access private keys or messages
    \item The platform resists government pressure for backdoors
    \item Security vulnerabilities are promptly identified and fixed
    \item User data is not monetized or shared with third parties
\end{itemize}

These trust assumptions have been repeatedly violated in practice, demonstrating the need for trustless alternatives.

\subsection{Legal and Regulatory Pressures}

Centralized platforms face significant legal challenges:
\begin{itemize}
    \item Government requests for user data and message access
    \item Requirements to implement lawful interception capabilities
    \item Pressure to scan messages for illegal content
    \item Compliance with data localization laws
    \item Liability for user-generated content
\end{itemize}

\subsection{Technical Vulnerabilities}

Centralized architectures create technical vulnerabilities:
\begin{itemize}
    \item Server-side vulnerabilities can expose all user data
    \item Compromised administrator accounts grant broad access
    \item Database breaches can leak encrypted keys and metadata
    \item Man-in-the-middle attacks during key exchange
    \item Client-side vulnerabilities in official applications
\end{itemize}

\section{Motivation for a Zero-Trust Approach}

\subsection{Principles of Zero-Trust Architecture}

Zero-trust security assumes that no component of the system should be inherently trusted. Applied to messaging systems, this means:

\begin{itemize}
    \item \textbf{Never Trust, Always Verify:} Every operation requires cryptographic proof
    \item \textbf{Least Privilege:} Components have minimal necessary permissions
    \item \textbf{Assume Breach:} Design assumes attackers have compromised the server
    \item \textbf{Defense in Depth:} Multiple independent security layers
\end{itemize}

\subsection{Benefits of Zero-Trust Messaging}

A zero-trust messaging system provides:

\begin{itemize}
    \item \textbf{Mathematical Privacy Guarantees:} Security based on cryptographic proofs, not trust
    \item \textbf{Reduced Attack Surface:} Server compromise does not expose message content
    \item \textbf{Regulatory Compliance:} Service provider cannot comply with content access requests
    \item \textbf{User Confidence:} Verifiable security properties increase user trust
    \item \textbf{Transparency:} Open-source implementation enables independent security audits
\end{itemize}

\subsection{Challenges in Zero-Trust Implementation}

Implementing zero-trust messaging introduces challenges:

\begin{itemize}
    \item \textbf{Key Management Complexity:} Users must manage their own cryptographic keys
    \item \textbf{Recovery Mechanisms:} Lost keys mean permanently lost messages
    \item \textbf{User Experience:} Cryptographic operations must be transparent
    \item \textbf{Performance Overhead:} Client-side encryption adds computational cost
    \item \textbf{Feature Limitations:} Some features (like server-side search) become impossible
\end{itemize}

\section{Positioning of This Project}

VoidLink adopts a hybrid approach that balances security, usability, and practicality:

\subsection{Architectural Positioning}

VoidLink uses a centralized server architecture with zero-trust security properties:
\begin{itemize}
    \item Centralized server for message routing and storage (usability)
    \item All cryptographic operations performed client-side (security)
    \item Server handles only opaque encrypted payloads (zero-trust)
    \item Two-layer authentication separates account and crypto operations
\end{itemize}

\subsection{Security Positioning}

VoidLink provides strong security guarantees while acknowledging limitations:

\textbf{Protected Against:}
\begin{itemize}
    \item Server operator accessing message content
    \item Database breaches exposing plaintext messages
    \item Network eavesdropping on message content
    \item Unauthorized message injection or modification
\end{itemize}

\textbf{Not Protected Against:}
\begin{itemize}
    \item Compromised client devices or malware
    \item Physical access to unlocked devices
    \item Coercion to decrypt messages
    \item Metadata analysis (who communicates with whom)
\end{itemize}

\subsection{Feature Positioning}

VoidLink deliberately limits features to maintain security:

\textbf{Implemented Features:}
\begin{itemize}
    \item One-to-one encrypted messaging
    \item Contact request and acceptance workflow
    \item Real-time message delivery via WebSocket
    \item Offline message queuing
    \item Encrypted key backup for multi-device access
    \item User presence and typing indicators
\end{itemize}

\textbf{Deliberately Excluded Features:}
\begin{itemize}
    \item Group messaging (increases cryptographic complexity)
    \item Server-side message search (requires plaintext access)
    \item Message forwarding (complicates attribution)
    \item Read receipts visible to server (metadata leakage)
    \item Social features (profile pictures, status updates)
\end{itemize}

\subsection{Target Use Cases}

VoidLink is designed for scenarios where privacy is paramount:
\begin{itemize}
    \item Confidential business communications
    \item Journalist-source communications
    \item Healthcare provider-patient messaging
    \item Legal professional communications
    \item Personal communications requiring strong privacy
\end{itemize}

The system prioritizes security and privacy over convenience, making it suitable for users who understand and accept the trade-offs inherent in zero-trust architectures.
\end{document}
